\section{Глава 11. *Методы с разделенной стратегией и аппроксимацией}

\subsection{Упражнение 11.1(с. 307)}
Преобразуйте уравнение (7.9) $n$-шагового TD-алгоритма с разделённой стратегией в
полуградиентную форму. Приведите также определения дохода в эпизодическом и
непрерывном случаях.
\begin{equation}
    V_{t+n}(S_t) \doteq V_{t+n-1}(S_t) + \alpha\,\rho_{t+1:t+n-1}
    \big[G_{t:t+n} - V_{t+n-1}(S_t)\big], \qquad 0 \le t < T,
    \tag{7.9}
\end{equation}
\begin{solution}
    По тому же рецепту, что и в (11.2): обновление массива $V$ заменяем
    обновлением весов $\mathbf w$, а приращение проецируем на $\nabla\hat v$.
    Обновление весов принимает вид
    \begin{equation*}
        \mathbf w_{t+n} \doteq \mathbf w_{t+n-1} + \alpha\,\rho_{t+1:t+n-1}
        \big[G_{t:t+n} - \hat v(S_t,\mathbf w_{t+n-1})\big]\nabla\hat v(S_t,\mathbf w_{t+n-1}),
        \quad 0\le t<T,
    \end{equation*}
    где коэффициент выборки по значимости
    \begin{equation*}
        \rho_{t+1:t+n-1} = \prod_{i=t+1}^{\min(t+n-1,\,T-1)} \frac{\pi(A_i\mid S_i)}{b(A_i\mid S_i)}.
    \end{equation*}

    Доход в эпизодическом случае (с обесцениванием):
    \begin{equation*}
        G_{t:t+n} \doteq R_{t+1} + \gamma R_{t+2} + \dots + \gamma^{n-1} R_{t+n}
        + \gamma^n \hat v(S_{t+n},\mathbf w_{t+n-1}), \qquad t+n < T,
    \end{equation*}
    и $G_{t:t+n} \doteq G_t$, если $t+n\ge T$.

    Доход в непрерывном случае (среднее вознаграждение, без обесценивания):
    \begin{equation*}
        G_{t:t+n} \doteq \big(R_{t+1}-\bar R_{t+n-1}\big) + \big(R_{t+2}-\bar R_{t+n-1}\big)
        + \dots + \big(R_{t+n}-\bar R_{t+n-1}\big) + \hat v(S_{t+n},\mathbf w_{t+n-1}),
    \end{equation*}
    то есть сумма дифференциальных наград за $n$ шагов плюс бутстрап $\hat
    v(S_{t+n})$, без $\gamma$.
\end{solution}

\subsection{Упражнение 11.2(с. 307)}
Преобразуйте уравнения (7.11) и (7.18) $n$-шагового $Q(\sigma)$ в полуградиентную
форму. Приведите также определения, охватывающие как эпизодический, так и
непрерывный случай.
\begin{equation}
    Q_{t+n}(S_t,A_t) \doteq Q_{t+n-1}(S_t,A_t)
    + \alpha\,\rho_{t+1:t+n}\big[G_{t:t+n} - Q_{t+n-1}(S_t,A_t)\big]
    \tag{7.11}
\end{equation}
\begin{equation}
    \begin{aligned}
    G_{t:h} \doteq{}& R_{t+1}
    + \gamma\big(\sigma_{t+1}\rho_{t+1} + (1-\sigma_{t+1})\pi(A_{t+1}\mid S_{t+1})\big)
      \big(G_{t+1:h} - Q_{h-1}(S_{t+1},A_{t+1})\big)\\
    &+ \gamma\bar V_{h-1}(S_{t+1}).
    \end{aligned}
    \tag{7.18}
\end{equation}
\begin{solution}
    Переводим в полуградиентную форму (массив $Q\to$ веса $\mathbf w$, приращение
    на $\nabla\hat q$). В $Q(\sigma)$ вся выборка по значимости сидит \emph{внутри}
    рекурсии дохода через коэффициент $\sigma_k\rho_k+(1-\sigma_k)\pi$, поэтому во
    внешнем обновлении весов множителя $\rho$ нет (иначе получилась бы двойная
    коррекция).

    Обновление весов:
    \begin{equation*}
        \mathbf w_{t+n} \doteq \mathbf w_{t+n-1}
        + \alpha\big[G_{t:t+n} - \hat q(S_t,A_t,\mathbf w_{t+n-1})\big]
        \nabla\hat q(S_t,A_t,\mathbf w_{t+n-1}).
    \end{equation*}

    Доход, эпизодический случай (с обесцениванием):
    \begin{equation*}
        \begin{aligned}
        G_{t:h} \doteq{}& R_{t+1}
        + \gamma\big(\sigma_{t+1}\rho_{t+1} + (1-\sigma_{t+1})\pi(A_{t+1}\mid S_{t+1})\big)
          \big(G_{t+1:h} - \hat q(S_{t+1},A_{t+1},\mathbf w_{t+n-1})\big)\\
        &+ \gamma\bar V_{h-1}(S_{t+1}),
        \end{aligned}
    \end{equation*}
    где ожидаемая ценность выражена через аппроксимацию (аналог (7.8)):
    \begin{equation*}
        \bar V_{h-1}(S_{t+1}) = \sum_a \pi(a\mid S_{t+1})\,\hat q(S_{t+1},a,\mathbf w_{t+n-1}).
    \end{equation*}

    Доход, непрерывный случай (среднее вознаграждение, без обесценивания): убираем
    $\gamma$ во всех местах и заменяем $R_{t+1}$ на $R_{t+1}-\bar R$:
    \begin{equation*}
        \begin{aligned}
        G_{t:h} \doteq{}& \big(R_{t+1}-\bar R\big)
        + \big(\sigma_{t+1}\rho_{t+1} + (1-\sigma_{t+1})\pi(A_{t+1}\mid S_{t+1})\big)
          \big(G_{t+1:h} - \hat q(S_{t+1},A_{t+1},\mathbf w_{t+n-1})\big)\\
        &+ \bar V_{h-1}(S_{t+1}).
        \end{aligned}
    \end{equation*}

    Рекурсия заканчивается на $G_{h:h}\doteq \hat q(S_h,A_h,\mathbf w_{t+n-1})$ при
    $h<T$ или на $G_{T-1:T}\doteq R_T$ при $h=T$ (здесь $h=t+n$).

    Две замены табличного на аппроксимацию: $Q_{h-1}(S_{t+1},A_{t+1})\to\hat
    q(S_{t+1},A_{t+1},\mathbf w)$ и $\bar V_{h-1}(S_{t+1})\to\sum_a\pi(a\mid
    S_{t+1})\hat q(S_{t+1},a,\mathbf w)$ (ожидание по $\pi$, а не одно действие).
    Коэффициент $\sigma_k\rho_k+(1-\sigma_k)\pi$ интерполирует между сэмплированием
    ($\sigma=1$: чистый importance sampling, сворачивается в наружное
    $\rho_{t+1:t+n}$ из (7.11)) и ожиданием ($\sigma=0$: tree-backup, $\rho$
    исчезает).
\end{solution}

\subsection{Упражнение 11.3(с. 311)}
Примените одношаговое полуградиентное Q-обучение к контрпримеру Бэрда и эмпирически покажите, что веса в этом случае расходятся.

\begin{solution}
    \textbf{Код смотри в отдельном файле репозитория.}
    Контрпример Бэрда: 7 состояний (6 верхних и 7-е), 2 действия (solid ведёт в
    7-е состояние, dashed --- в случайное верхнее), награда всюду $0$,
    $\gamma=0.99$. Поведенческая стратегия $b$ берёт dashed с вероятностью $6/7$ и
    solid с $1/7$ (даёт равномерное посещение). Линейная аппроксимация на восьми
    весах со структурой признаков Бэрда: ценность solid равна $w_7+2w_8$, ценность
    dashed из верхнего состояния $i$ равна $2w_i+w_8$; общий вес $w_8$ входит во
    все состояния. Обновление --- одношаговое Q-обучение $\mathbf w \leftarrow
    \mathbf w + \alpha\big[\gamma\max_a \hat q(S',a) - \hat q(S,A)\big]\nabla\hat
    q(S,A)$ (без коэффициента выборки по значимости: цель строится по жадному
    действию напрямую). Начальные веса $\mathbf w=(1,1,1,1,1,1,10,1)$, $\alpha=0.01$.

    \begin{figure}[H]
        \centering
        \includegraphics[width=0.8\textwidth]{fig_sol_11_3.png}
        \caption*{Эволюция весов при Q-обучении на контрпримере Бэрда.}
    \end{figure}

    Веса расходятся, как на рис.~11.2: $w_8$ взлетает (за $1000$ шагов до
    $\approx 190$, а к $5000$ шагам --- порядка $10^5$), $w_1..w_6$ растут, $w_7$
    остаётся сравнительно ограниченным. Рост неограниченный и не зависит от
    $\alpha>0$.

    Механизм --- смертельная триада. Из-за начального $w_7=10$ ценность 7-го
    состояния велика, поэтому $\max_a$ устойчиво указывает на solid, и цель всех
    обновлений бутстрапит на $\hat q(\cdot,\text{solid})=w_7+2w_8$. Обновления идут
    по off-policy распределению $b$ (в $6/7$ случаев --- на верхних состояниях, где
    dashed поднимает общий $w_8$). Но $w_8$ входит в саму цель с коэффициентом $2$,
    поэтому цель убегает вверх быстрее, чем за ней подтягиваются оценки: положительная
    обратная связь. Компенсирующих обновлений (которые вернули бы $\hat q(7)$ на
    место) поведенческая стратегия почти не даёт --- отсюда расходимость. Убрать
    любой элемент триады (аппроксимацию, бутстрэппинг или off-policy распределение)
    --- и веса сходятся.
\end{solution}
\notebook{11_03}{11.3}

\subsection{*Упражнение 11.4(с. 324)}
Докажите тождество (11.24):
\begin{equation}
    \overline{RE}(\mathbf w) = \mathbb E\big[(G_t - \hat v(S_t,\mathbf w))^2\big]
    = \overline{VE}(\mathbf w) + \mathbb E\big[(G_t - v_\pi(S_t))^2\big].
    \tag{11.24}
\end{equation}

\begin{solution}
    Обозначения: $\mathbb E[\,\cdot\mid S_t=s]$ --- усреднение по случайности дохода
    $G_t$ при фиксированном $s$; безусловное матожидание раскрывается по закону
    полного матожидания $\mathbb E[X]=\sum_s\mu(s)\,\mathbb E[X\mid S_t=s]$.
    Пользуемся тем, что при фиксированном $s$ величины $\hat v(s,\mathbf w)$ и
    $v_\pi(s)$ --- константы, а $\mathbb E[G_t\mid S_t=s]=v_\pi(s)$ по определению
    истинной ценности.

    Раскрываем $\overline{RE}$ в условное матожидание:
    \begin{equation*}
        \overline{RE}(\mathbf w) = \sum_s \mu(s)\,
        \mathbb E\big[(G_t - \hat v(s,\mathbf w))^2 \mid S_t=s\big].
    \end{equation*}
    Добавляем и вычитаем $v_\pi(s)$ и группируем:
    \begin{equation*}
        G_t - \hat v(s,\mathbf w) = \underbrace{\big(G_t - v_\pi(s)\big)}_{A}
        + \underbrace{\big(v_\pi(s) - \hat v(s,\mathbf w)\big)}_{B},
    \end{equation*}
    где при фиксированном $s$ величина $A$ случайна (зависит от $G_t$), а $B$ ---
    константа. Раскрываем квадрат под условным матожиданием:
    \begin{equation*}
        \mathbb E\big[(A+B)^2 \mid s\big]
        = \mathbb E[A^2\mid s] + 2\,\mathbb E[AB\mid s] + \mathbb E[B^2\mid s].
    \end{equation*}

    Разбираем члены. Так как $B$ --- константа при данном $s$:
    \begin{align*}
        \mathbb E[B^2\mid s] &= \big(v_\pi(s) - \hat v(s,\mathbf w)\big)^2,\\
        2\,\mathbb E[AB\mid s] &= 2B\,\mathbb E[A\mid s]
        = 2B\big(\underbrace{\mathbb E[G_t\mid S_t=s]}_{=\,v_\pi(s)} - v_\pi(s)\big) = 0.
    \end{align*}
    Перекрёстный член зануляется не из-за независимости, а потому что $B$
    выносится как константа, а $\mathbb E[A\mid s]=0$. Остаётся
    \begin{equation*}
        \mathbb E\big[(G_t - \hat v(s,\mathbf w))^2 \mid s\big]
        = \big(v_\pi(s) - \hat v(s,\mathbf w)\big)^2
        + \mathbb E\big[(G_t - v_\pi(s))^2 \mid s\big].
    \end{equation*}

    Возвращаем внешнее усреднение $\sum_s\mu(s)$:
    \begin{equation*}
        \overline{RE}(\mathbf w)
        = \underbrace{\sum_s \mu(s)\big(v_\pi(s) - \hat v(s,\mathbf w)\big)^2}_{\overline{VE}(\mathbf w)}
        + \underbrace{\sum_s \mu(s)\,\mathbb E\big[(G_t - v_\pi(s))^2\mid S_t=s\big]}_{\mathbb E[(G_t - v_\pi(S_t))^2]}.
    \end{equation*}
\end{solution}
